%% 
% 维护规则
% 新增数学推导前，先检查这里是否已有同名符号；若冲突，优先改章内符号，保证全书一致。
% 需要矩阵/向量排版规则（粗体、黑板体、希腊字母）时，在表头下面加一条“排版约定”注释：
%- 向量/矩阵使用粗体罗马体：$\mathbf{x}, \mathbf{W}$；
%-集合/域使用黑板体：$\mathbb{R}, \mathbb{N}$；
%-标量使用斜体：$x, y, \lambda$。
% 本表同样用 %% 
% 维护规则
% 新增数学推导前，先检查这里是否已有同名符号；若冲突，优先改章内符号，保证全书一致。
% 需要矩阵/向量排版规则（粗体、黑板体、希腊字母）时，在表头下面加一条“排版约定”注释：
%- 向量/矩阵使用粗体罗马体：$\mathbf{x}, \mathbf{W}$；
%-集合/域使用黑板体：$\mathbb{R}, \mathbb{N}$；
%-标量使用斜体：$x, y, \lambda$。
% 本表同样用 %% 
% 维护规则
% 新增数学推导前，先检查这里是否已有同名符号；若冲突，优先改章内符号，保证全书一致。
% 需要矩阵/向量排版规则（粗体、黑板体、希腊字母）时，在表头下面加一条“排版约定”注释：
%- 向量/矩阵使用粗体罗马体：$\mathbf{x}, \mathbf{W}$；
%-集合/域使用黑板体：$\mathbb{R}, \mathbb{N}$；
%-标量使用斜体：$x, y, \lambda$。
% 本表同样用 %% 
% 维护规则
% 新增数学推导前，先检查这里是否已有同名符号；若冲突，优先改章内符号，保证全书一致。
% 需要矩阵/向量排版规则（粗体、黑板体、希腊字母）时，在表头下面加一条“排版约定”注释：
%- 向量/矩阵使用粗体罗马体：$\mathbf{x}, \mathbf{W}$；
%-集合/域使用黑板体：$\mathbb{R}, \mathbb{N}$；
%-标量使用斜体：$x, y, \lambda$。
% 本表同样用 \input{symbol_table.tex} 引入（通常放在前言或附录）

% symbol_table.tex 符号表（集中维护）
\begin{table}[htbp]
\centering
\caption{通用符号约定}
\begin{tabular}{@{}p{2.2cm}p{7.2cm}p{3.0cm}p{3.0cm}p{2.2cm}@{}}
\toprule
符号 & 含义 & 域/形状/单位 & 默认约定 & 首次出现 \\
\midrule
$\mathbf{x}$ & 输入特征向量 & $\mathbb{R}^{d}$ & 列向量，粗体表向量 & 第7章 \\
$y$ & 监督标签 & $\{0,1\}$ 或 $\mathbb{R}$ & 分类/回归上下文限定 & 第7章 \\
$\hat{y}$ & 模型预测 & 同 $y$ & argmax 或回归输出 & 第7章 \\
$\mathbf{W}$ & 权重矩阵 & $\mathbb{R}^{n\times d}$ & 行对应单元，列对应特征 & 第7章 \\
$b$ & 偏置 & $\mathbb{R}$ & 标量 & 第7章 \\
$\mathcal{L}$ & 损失函数 & — & 统一写作“损失（loss）” & 第7章 \\
$\lambda$ & 正则系数 & $\mathbb{R}_{\ge 0}$ & L2 默认为权重衰减 & 第7章 \\
$\eta$ & 学习率 & $\mathbb{R}_{>0}$ & SGD/Adam上下文说明 & 第7章 \\
$\sigma(\cdot)$ & 激活函数 & — & 章内注明具体形式 & 第7章 \\
$\epsilon$ & 扰动幅度 & $\mathbb{R}_{\ge 0}$ & 对抗样本/数值稳定 & 第10章 \\
\bottomrule
\end{tabular}
\end{table} 引入（通常放在前言或附录）

% symbol_table.tex 符号表（集中维护）
\begin{table}[htbp]
\centering
\caption{通用符号约定}
\begin{tabular}{@{}p{2.2cm}p{7.2cm}p{3.0cm}p{3.0cm}p{2.2cm}@{}}
\toprule
符号 & 含义 & 域/形状/单位 & 默认约定 & 首次出现 \\
\midrule
$\mathbf{x}$ & 输入特征向量 & $\mathbb{R}^{d}$ & 列向量，粗体表向量 & 第7章 \\
$y$ & 监督标签 & $\{0,1\}$ 或 $\mathbb{R}$ & 分类/回归上下文限定 & 第7章 \\
$\hat{y}$ & 模型预测 & 同 $y$ & argmax 或回归输出 & 第7章 \\
$\mathbf{W}$ & 权重矩阵 & $\mathbb{R}^{n\times d}$ & 行对应单元，列对应特征 & 第7章 \\
$b$ & 偏置 & $\mathbb{R}$ & 标量 & 第7章 \\
$\mathcal{L}$ & 损失函数 & — & 统一写作“损失（loss）” & 第7章 \\
$\lambda$ & 正则系数 & $\mathbb{R}_{\ge 0}$ & L2 默认为权重衰减 & 第7章 \\
$\eta$ & 学习率 & $\mathbb{R}_{>0}$ & SGD/Adam上下文说明 & 第7章 \\
$\sigma(\cdot)$ & 激活函数 & — & 章内注明具体形式 & 第7章 \\
$\epsilon$ & 扰动幅度 & $\mathbb{R}_{\ge 0}$ & 对抗样本/数值稳定 & 第10章 \\
\bottomrule
\end{tabular}
\end{table} 引入（通常放在前言或附录）

% symbol_table.tex 符号表（集中维护）
\begin{table}[htbp]
\centering
\caption{通用符号约定}
\begin{tabular}{@{}p{2.2cm}p{7.2cm}p{3.0cm}p{3.0cm}p{2.2cm}@{}}
\toprule
符号 & 含义 & 域/形状/单位 & 默认约定 & 首次出现 \\
\midrule
$\mathbf{x}$ & 输入特征向量 & $\mathbb{R}^{d}$ & 列向量，粗体表向量 & 第7章 \\
$y$ & 监督标签 & $\{0,1\}$ 或 $\mathbb{R}$ & 分类/回归上下文限定 & 第7章 \\
$\hat{y}$ & 模型预测 & 同 $y$ & argmax 或回归输出 & 第7章 \\
$\mathbf{W}$ & 权重矩阵 & $\mathbb{R}^{n\times d}$ & 行对应单元，列对应特征 & 第7章 \\
$b$ & 偏置 & $\mathbb{R}$ & 标量 & 第7章 \\
$\mathcal{L}$ & 损失函数 & — & 统一写作“损失（loss）” & 第7章 \\
$\lambda$ & 正则系数 & $\mathbb{R}_{\ge 0}$ & L2 默认为权重衰减 & 第7章 \\
$\eta$ & 学习率 & $\mathbb{R}_{>0}$ & SGD/Adam上下文说明 & 第7章 \\
$\sigma(\cdot)$ & 激活函数 & — & 章内注明具体形式 & 第7章 \\
$\epsilon$ & 扰动幅度 & $\mathbb{R}_{\ge 0}$ & 对抗样本/数值稳定 & 第10章 \\
\bottomrule
\end{tabular}
\end{table} 引入（通常放在前言或附录）

% symbol_table.tex 符号表（集中维护）
\begin{table}[htbp]
\centering
\caption{通用符号约定}
\begin{tabular}{@{}p{2.2cm}p{7.2cm}p{3.0cm}p{3.0cm}p{2.2cm}@{}}
\toprule
符号 & 含义 & 域/形状/单位 & 默认约定 & 首次出现 \\
\midrule
$\mathbf{x}$ & 输入特征向量 & $\mathbb{R}^{d}$ & 列向量，粗体表向量 & 第7章 \\
$y$ & 监督标签 & $\{0,1\}$ 或 $\mathbb{R}$ & 分类/回归上下文限定 & 第7章 \\
$\hat{y}$ & 模型预测 & 同 $y$ & argmax 或回归输出 & 第7章 \\
$\mathbf{W}$ & 权重矩阵 & $\mathbb{R}^{n\times d}$ & 行对应单元，列对应特征 & 第7章 \\
$b$ & 偏置 & $\mathbb{R}$ & 标量 & 第7章 \\
$\mathcal{L}$ & 损失函数 & — & 统一写作“损失（loss）” & 第7章 \\
$\lambda$ & 正则系数 & $\mathbb{R}_{\ge 0}$ & L2 默认为权重衰减 & 第7章 \\
$\eta$ & 学习率 & $\mathbb{R}_{>0}$ & SGD/Adam上下文说明 & 第7章 \\
$\sigma(\cdot)$ & 激活函数 & — & 章内注明具体形式 & 第7章 \\
$\epsilon$ & 扰动幅度 & $\mathbb{R}_{\ge 0}$ & 对抗样本/数值稳定 & 第10章 \\
\bottomrule
\end{tabular}
\end{table}